\chapter{Conclusions and Future Work}
\label{ch:conclusions}

\section{Summary of Contributions}
\label{sec:conclusions_summary}

% TODO: Three-paragraph summary.
%
% 1. We built a real-time voice-to-instrument resynthesis system using
%    BRAVE deployed in Pure Data with sub-10 ms latency. The drums
%    prototype (drums_v1 ep755, later drums_v2 if successful) responds
%    audibly to voice and beatbox input.
%
% 2. We trained a working guitar autoencoder (guitar_v6 with the official
%    RAVE v2 architecture) which reconstructs guitar from guitar input
%    with recognisable pluck-like character.
%
% 3. We instrumented six guitar training iterations with quantitative
%    diagnostics and produced a failure taxonomy that identifies and
%    distinguishes (a) posterior collapse, (b) partial latent collapse
%    from oversized representations, (c) GAN discriminator dominance,
%    (d) GAN discriminator collapse, and (e) decoder mode collapse on
%    out-of-distribution input. This taxonomy is, to our knowledge, the
%    most detailed published account of small-dataset RAVE failure modes.

\section{Limitations}
\label{sec:conclusions_limitations}

\subsection{Voice-to-Guitar Did Not Succeed}
% TODO: Honest framing. The encoder is not the bottleneck (it
% differentiates voice inputs in latent space). The decoder is not
% broken (it reconstructs guitar from guitar input). The mismatch is
% that voice and guitar live in disjoint regions of the learned latent
% manifold, and the decoder is not trained to produce meaningful output
% from voice-encoded latents. Two well-known workarounds were not
% pursued here due to time constraints (see Future Work).

\subsection{Noise-Floor Heterogeneity}
% TODO: The 35.8 dB spread across recording methods means the model
% encodes recording-method as a salient feature of its latent space.
% A truly clean training corpus would restrict to a single recording
% method.

\subsection{Phase 2 GAN Instability}
% TODO: Even with the v2 architecture's safeguards, Phase 2 dynamics
% are not monotonically improving -- the model peak around mid-Phase 2
% can degrade by the end. We pick the best checkpoint manually rather
% than trusting the final state.

\section{Future Work}
\label{sec:conclusions_future}

\subsection{Voice Pre-Processing for In-Distribution Mapping}
% TODO: Real-time pitch extraction (e.g., CREPE, fiddle~, sigmund~) plus
% harmonic resynthesis would let us synthesise a guitar-like signal
% from the voice's pitch and energy envelope, which would then fall
% inside the encoder's training distribution. This is the DDSP-style
% approach applied as a pre-processing front-end.

\subsection{Solo-Only Guitar Training}
% TODO: Restricting to GuitarSet's 180 _solo_mic files (monophonic
% playing) plus selected monophonic content from IDMT-SMT-Guitar would
% produce a model whose latent space matches voice's monophonic nature.
% Smaller but cleaner corpus.

\subsection{Larger and More Consistent Datasets}
% TODO: A single-source, large guitar corpus (e.g., GOAT, GAPS, or a
% commercially licensed library) would eliminate noise-floor
% heterogeneity. The BRAVE paper's saxophone success was on
% single-source data.

\subsection{Alternative Architectures}
% TODO: DDSP, SoundStream-style discrete-token models, or diffusion-based
% audio generation (e.g., DiffSinger / DiffAudio variants) might better
% handle voice-to-instrument mapping.

\section{Closing Remarks}
\label{sec:conclusions_closing}

% TODO: One-paragraph wrap. Even when an ML system does not solve the
% headline task, careful instrumentation of its failures is a useful
% scientific output. The diagnostic methodology developed for this thesis
% would apply to any RAVE-family training and would surface the same
% failure modes faster than the trial-and-error path we travelled.
